\chapter{CONCLUSIONS AND FUTURE WORK}
\label{ch:conclusions}

\section{Introduction}
The main goal in this dissertation has been to create the conceptual framework for a Localized High-Performance Order Management System that is designed for use with the unpredictable and very-high volume nature of temporary cultural events and festivals. A variety of factors exist which can limit or cause failure within a traditional cloud-centric, synchronous architecture. The majority of such systems are not able to function at all in unreliable Ad-Hoc networked environments, they also typically cannot handle large amounts of patron activity simultaneously as well as do not have sufficient hardware resources. To address each of these limitations as they relate to the potential for system failure, the researcher created an asynchronous, event-driven environment. In this final chapter we will summarize the major technological accomplishments from the project and show how the results of our empirical analysis aligns with the initial goals of the project. We also outline some possible areas where additional academic research and business development could occur.

\section{Synthesis of Technical Accomplishments}
This thesis demonstrates a fully realized software engineering life cycle from architecture through implementation. The MERN (MongoDB, Express.js, React, Node.js) technology stack \cite{banker2011mongodb, express, freeman2019pro, tilkov2010node} was heavily modified for localized, off-line first topologies during the architectural design and implementations. In order to replace the traditional poll based request/response HTTP protocol using the full duplex nature of WebSockets \cite{fette2011websocket, pimentel2012communicating} the system created bi-directional communications between each terminal. Additionally, the user interface was built around optimizations for fast rendering. The use of the virtual dom along with a bespoke dual-mode design system ('Eclipse UI'), supporting seamless light/dark theme switching while maintaining consistent high-contrast, color-coded interactive components across all operational modules \cite{freeman2019pro} provided high readability and 60 frames per second performance while continuously rendering large amounts of real time data.

During the development of the back-end, atomicity of states and strong recover mechanisms were prioritized. Using the centralized event handler model allowed the avoidance of relational database locking issues by maintaining state and operations entirely within memory until an immutable transaction could be committed to the MongoDB database. State-less JSON Web Tokens \cite{ferraiolo1992rbac} were used to implement role-based access control as well as ensure strict data integrity within a zero trust local area network \cite{nist2020zero}. Finally, the agile/scrum method \cite{scrum2020guide} was used to iteratively improve over several sprint cycles and continually increase the overall resilience against various edge case scenarios found in both lab testing and actual production deployments.

\section{Mapping to Research Objectives}
The success criteria that were specified as part of the research objectives, which are outlined in Chapter 1 Section 1.3, provide the tangible basis to evaluate the success of this proposed system. A direct comparison is provided in Table~\ref{tab:objectives_map}, between each of the primary research objectives that were identified, and the evidence-based data collected from the primary research activities conducted during the evaluation phase as described in Chapter 5. All of the primary research objectives have successfully been met.

\begin{table}[htbp]
\centering
\small
\begin{tabular}{|p{4.5cm}|p{5cm}|p{2.2cm}|}
\hline
\textbf{Research Objective (Ch.~1)} & \textbf{Empirical Evidence (Ch.~5)} & \textbf{Status} \\ \hline
Support $\geq$50 concurrent client devices & Sustained 500 concurrent connections with $<$50ms P95 latency (Table~\ref{tab:comparison_results}) & \textbf{Exceeded} \\ \hline
Process $\geq$1,000 transactions per hour & 5,247 req/sec sustained throughput (Table~\ref{tab:stress_results}), equating to $\sim$18.9M/hour & \textbf{Exceeded} \\ \hline
Operate on commodity hardware ($<$2GB RAM) & 50 KB/client; 108.5 MB total footprint; 38\% peak CPU & \textbf{Fulfilled} \\ \hline
Eliminate phantom orders; absolute data consistency & 0\% error rate across 1.58M samples; atomic FIFO queue (Section 5.4.3) & \textbf{Fulfilled} \\ \hline
Real-time reactive dashboards for management & Dashboard deployed throughout Cultural Week field study (Figure~\ref{fig:admin_heatmap}) & \textbf{Fulfilled} \\ \hline
Resilient operation under network instability & Sub-310ms RTO after server failure; offline buffer recovery (Section 5.4.4) & \textbf{Fulfilled} \\ \hline
\end{tabular}
\caption{Systematic mapping of the original research objectives to the empirical evidence gathered during the evaluation phase, confirming fulfilment of all stated goals.}
\label{tab:objectives_map}
\end{table}

\section{Primary Conclusions and Research Findings}
The empirically based evaluation performed at the end of this research verified the hypotheses formulated at the beginning of the project and made possible many important conclusions about architectures for distributed Order Management Systems.

First, the base line comparison studies, rigorously contextualized using the classic queueing theory \cite{kleinrock1975} and quantitatively determined through Little's Law \cite{little1961}, definitively showed the advantages in terms of scalability in the use of an asynchronous model with event driven behavior compared to a traditional model that uses synchronous thread blocking. For example, the implementation was able to process hundreds of clients concurrently with a reduction in processing time of 62\%, and a reduction in memory allocations per client of 97\% (from 1.5 MB to 50 KB) when compared to traditional implementations. The implementation also demonstrated how resilient it was to spiky latency events common in the case where server thread pools become saturated the average response time remained less than 50 milliseconds, while the average response time for the baseline synchronous implementation was greater than 800 milliseconds.

Secondly, the evaluation supported the conclusion that moving from RESTful API polling to raw WebSocket data framing \cite{fette2011websocket} results in a decrease of 80.7\% in network overhead per transaction (from 1450 bytes to 280 bytes), resulting in a significant preservation of wireless spectrum. This resulted in a reduction in packet collisions on the local area network used by the mobile device to communicate with the central service \cite{ibm2019cloud}.

Thirdly, the evaluations of resilience and fault tolerance using chaos engineering techniques \cite{basiri2016chaos} showed that the system was capable of maintaining its transactional state even after having been subjected to extreme levels of network fragmentation. Using locally stored buffers to isolate outgoing operations until connections could be reestablished, and then automatically synchronizing operations once connectivity had been restored eliminated two major issues related to distributed systems: phantom orders and financial dis-synchronizations. Phantom orders occur due to the fact that some systems do not properly track changes to order status, thereby leading to duplicate orders being processed. Financial dis-synchronizations occur because the system does not have a proper mechanism to ensure that both sides of a transaction have reached agreement before updating their records. Both of these problems can lead to serious consequences if they are allowed to go unchecked. Examples include situations where funds may be deducted from one account but not credited to another. These types of problems are well known in the literature as examples of ``distributed computing paradoxes,'' such as the Two Generals Problem \cite{brewer2012cap}. The Recovery Time Objective (RTO) is defined as the amount of time it takes for a system to recover from a failure and resume normal operation. A RTO of 310 milliseconds indicates that this system is capable of recovering from most types of failures autonomously, without requiring manual administrative intervention.

Finally, the evaluation campaign focused on validating the security attributes of the layered defense-in-depth architecture demonstrated that it has successfully defended against all of the attack vectors identified in the OWASP Top 10:2025 list \cite{owasp2025} specifically against SQL injection attacks, cross site scripting attacks, attacks related to broken access controls, and user enumeration attacks.

In summary, the closed lab-based measurement results were significantly reinforced by large scale operational deployment tests. During testing, we maintained 694 transactions totaling a financial turnover exceeding \texteuro 22,000, and did so without experiencing any failures under high volume bursts of traffic. In essence, our physical event demonstration provided clear evidence that our solutions would be highly applicable and useful for large scale industrial and commercial applications.

\section{Limitations of the Current Work}
In spite of the fact that the constructed system has successfully accomplished each one of the primary research goals, it is necessary for these limitations to be recognized to provide context regarding the scope and generalization of the results.

\textbf{Single-server topology.} The present deployment configuration relies upon a single central server instance, which will function as a backend. Although the results from the chaos testing demonstrated that the systems could rapidly recover automatically (RTO=310ms) there are no present capabilities within the system to utilize active/active fail-over or horizontal scaling between servers. If a catastrophic hardware failure occurs and there is not a standby server available, the entire system would need to be manually restored before availability would be returned to the users.

\textbf{Single event field data collection.} Telemetry data were collected from only one regional cultural festival at a time with a limited number of patrons and operations. While the 694 transaction data set provides extensive validation of system reliability, the findings have not been validated under substantially larger multi-day music festivals, geographically distributed multi-venue configurations, or operationally different culturally-based environments.

\textbf{Not PCI-DSS compliant credit card processing gateway.} There was also previously stated limitation of scope (Chapter 1, Section 1.4) as well as the original scope documentation. This indicates that the system does not contain a PCI-DSS compliance approved payment gateway. All financial transactions that occurred during the field deployment were handled through either cash, bank transfers, or external card terminal equipment that did not integrate into the WebSocket architecture.

\textbf{Synchronous comparative baseline method.} The baseline used for comparative purposes relied upon the Flask development server as opposed to production grade WSGI deployments as described in the threats to validity section of Chapter 5. The selection of the baseline was made intentionally in order to clearly demonstrate and identify the core architectural differences; however, the differential in actual performance may differ when comparing against an optimized multi-worker synchronous deployment.

\section{Future Work and Architectural Expansion}
While the developed Order Management System successfully fulfilled its primary objectives, the dynamic nature of distributed computing presents multiple avenues for future academic research and commercial expansion. The following directions are prioritized by estimated research impact and implementation feasibility.

\subsection{Near-Term: Predictive Analytics and Intelligent Inventory Management}
The effect of this integration effort is to incorporate predictive models of machine learning into the existing local server architecture so that the models may be used with historical velocities for every transaction and the actual temporal density (inter-arrival distribution) as found in Table~\ref{tab:intensity_analysis}, to predict the timing of when items at a given station will run out. These systems are capable of sending automatic notifications to management to replenish items located at specific stations based on the predicted times and therefore increase the service rate by even more than the original modifications. The initial training dataset for all of the models is comprised of the annotated data set from the field study containing 694 transactions.

\subsection{Medium-Term: Multi-Node Replication via CRDTs}
A major architectural improvement includes the use of CRDTs \cite{shapiro2011crdt} as an example of conflict-free replicated data types to enable simple multi-node master-master clustering. The current system is very reliable on one server at a time. However, by having the main database split over many synchronized servers spread out across a large area, the number of potential faults would be greatly increased and all the limitations of relying on a single server would be removed as described in section 6.5. Therefore, identifying the mathematical consensus protocols used to resolve conflicts that occur when merging the states from different local nodes which are not controlled by a common central arbiter specifically for counter based inventory updates and set based orders is a very difficult yet worthwhile research goal and will directly influence the tradeoffs of the CAP theorem identified in chapter 2 \cite{brewer2012cap}.

\subsection{Long-Term: NFC-Based Cashless Payment Integration}
The integration capabilities of the system can be extended to include Near Field Communication (NFC) protocol and Centralized Radio-Frequency Identification (RFID) closed loop payment systems. Integration of a Localized Cashless Wristband used in large global festival events directly with the Web Sockets Infrastructure will eliminate all outside Banking Gateway Requirements. In theory it should provide the greatest possible optimization to System Throughput as well as Consumer Satisfaction by reducing Check Out Time to under One Second Per Transaction. However, this is an extremely ambitious path that would require significant research into Financial Regulatory Compliance Frameworks as well as Secure Element Provisioning.

\section{Closing Statement}
The research in this dissertation has shown that the integration of modern asynchronous web technologies online/offline first deployment methodologies and event driven design principles create an architecture that is significantly superior to legacy cloud-centric systems in both volatile and very high intensity environment. The architecture that was developed through the elimination of the dependency on external internet gateways and the relocation of computing authority from remote servers to the local network, provided a zero latency synchronization capability, ensured complete operational independence; and provided a significant increase in hardware utilization.

As we move forward, this architectural model represents a significant change in how mission critical applications are designed for localized events. The future of high density transactional processing will not come from continually increasing the scale of centralize cloud-based resources to handle unpredictable surges of human activity, but instead will rely on taking back control and distributing the decision making authority to the operational edge. Autonomous event driven architectures utilizing local first design principles, coupled with peer to peer data structures that are fault tolerant, represent the ultimate level of digital sovereignty: it provides the certainty that regardless of the condition of the global internet, your application will continue to operate at the local level.

The research contributions outlined in this dissertation provide a robust, flexible base for realizing this vision, and demonstrate that all the necessary tools exist today to develop the next-generation of fully autonomous, failure-free local networks.
